% 1 pag (Max 2)
%Consiste en identificar los fenómenos, hechos o situaciones, que puestos en relación presentan incongruencia, obstáculos, desconocimiento o discrepancia y que constituyen el objeto de estudio

% Problema Cientifico 

La localización de las fuentes precisas en estudios de electroencefalógrafo (EEG) constituye una gran herramienta para el estudio no invasivo del cerebro humano. Sin embargo, aun con el desarrollo tecnológico de las ultimas decadas siguen existiendo desafíos tanto en un nivel técnico como metodológico esto hace que se limiten sus aplicaciones principalmente en contextos donde la precisión es esencial o en condiciones de recolección donde existen altos indices de ruido. La comunidad científica ha identificado que las técnicas más recientes como MNE, sLORETA y VARETA que están basadas en modelos inversos, tienden a favorecer fuentes ubicadas en regiones más superficiales del cerebro, esto ocasiona que se presenten dificultades notables para detectar fuentes en un nivel más profundo o distribuidas, lo que conlleva un sesgo sistemático en los resultados.

Adicionalmente, las métricas empleadas en estudios anteriores tienden a no evaluar de manera integral la calidad de la localización, mostrando amplias regiones de actividad sin llegar a ser precisas. Muchas se reducen a promedios o zonas de probabilidad las cuales no diferencian entre precisión espacial, resistencia al ruido o adaptabilidad entre individuos. Esto hace que sea más difícil la comparación objetiva entre métodos y la creación de procesos estandarizados. Ademas de esto se suma que las pruebas de rendimiento  se suelen realizar con conjuntos de datos reducidos, poco diversos o con controles experimentales insuficientes, esto ocasiona que se comprometa la fiabilidad estadística de los resultados. Es por esto que resulta complejo determinar si las fallas de los métodos actuales se deben a sus bases teóricas, a las condiciones de adquisición de datos o a los criterios de validación empleados.

Otro problema importante es la desconexión entre los modelos computacionales y su implementación real. Muchos métodos teóricos no consiguen aplicarse de manera eficiente a los datos reales, que contienen varios factores no considerados en la teoría como la variabilidad de los artefactos, variabilidad anatómica y diferencias individuales en la conductividad craneal. Estas diferencias entre la teoría y la práctica crean deficiencias para su paso a entornos de aplicación como lo son ambientes clínicos o neurotecnológicos, haciendo que la tecnología no sea completamente accesible. Si a esto sumamos la falta de herramientas para visualizar de manera sencilla así como para interpretar los mapas de activación, nos enfrentamos a un limitado uso en contextos interdisciplinarios.

Teniendo en cuenta estas limitantes las cuales han sido ampliamente identificadas en investigaciones recientes, se destaca la necesidad de metodologías más sólidas, flexibles y reproducibles. Sin embargo, el enfoque de las soluciones propuestas está en mejoras marginales de modelos ya existentes, sin detenerse a resolver de manera integral los factores que impactan la localización o sin proponer nuevas metodologías. Del mismo modo, la interpretación de manera sencilla de los resultados sigue siendo crucial para usuarios no expertos en procesamiento de señales de manera que sea más accesible para todos y de esta manera tener un mayor alcance.

